\documentclass[11pt]{article}
\usepackage[a4paper,left=2.6cm,right=2.6cm,top=2.6cm,bottom=2.8cm]{geometry}
\usepackage[T1]{fontenc}\usepackage[utf8]{inputenc}\usepackage{mathptmx}\usepackage{microtype}
\usepackage[hidelinks]{hyperref}
\setlength{\parindent}{0pt}\setlength{\parskip}{4pt}
\pagestyle{plain}
\begin{document}
\begin{center}
{\LARGE\bfseries The Price of a Design Parameter}\\[4pt]
{\large Bank Valuations in the Digital Euro Legislative Process}\\[12pt]
{\large Philip Kroos}\\[6pt]
October 2026
\end{center}
\vspace{6pt}

\textbf{Question.} The principal financial-stability concern surrounding a retail central bank digital currency is deposit flight from commercial banks. The digital euro addresses this risk through concrete design parameters, including a holding limit, rules for corporate holdings, fee caps and access requirements. This paper asks whether equity markets priced these design choices during the European legislative process, and whether the response was larger for firms with greater economic exposure.

\textbf{Data and identification.} I assemble 26 legislative and central-bank events between 2023 and 2026, of which 22 verified event dates enter the pre-specified estimation sample, and code their contemporaneously available design content separately for banks, payment providers and device gatekeepers. I combine these events with daily equity returns for 49 listed firms, pre-proposal EBA measures of bank deposit exposure, and the published record of 309 disclosed meetings between legislators and interest representatives. The main specification is a two-stage event study. Abnormal returns are first estimated from market-factor models; the second stage relates firm-event abnormal returns to exposure interacted with the coded direction of each event, with firm and event fixed effects. Identification therefore comes from cross-sectional differences in exposure on the same event day and from sign changes in the coded design content across events.

\textbf{Findings.} When every event is dated by the day its content became public, the relative gain of a euro-area bank one standard deviation more deposit-funded, on an event that limits the digital euro, is $-0.01$ percentage points (placebo-date $p = 0.91$). The pre-specified dating, which used the document date of a draft report already covered by the press, gives 0.05 ($p = 0.45$). Allowing for shocks common to all banks on an event day, the confidence bounds exclude effects about a quarter to a third the size of those documented for the early digital-euro announcements. Events that set or signal a design parameter show no detectable valuation effect. As a positive control, using predetermined exposure and strictly pre-event factor loadings, the same design finds a response of the documented size on the day of the October 2020 Eurosystem report, larger in absolute value than 94\% of the slopes on ordinary trading days.

\textbf{Interpretation and contribution.} The results suggest that markets reacted more strongly to early news about whether a digital euro would exist than to the later legislative calibration of its parameters. This is not because the parameters were politically unimportant: the negotiation record shows intense contestation, with the banking side the most frequent private counterparty. The paper contributes a market-based test of deposit-flight risk from CBDC design, an event-level coding of the digital euro legislative process, and a new public dataset on the political economy of that process.

\textbf{Keywords:} central bank digital currency; digital euro; financial stability; deposit flight; banking; event study; political economy of regulation.
\end{document}
